% =====================================================================
%  Lecture 14 -- 5G Networks: Introduction                 (Persian edition)
% =====================================================================
\documentclass[11pt,a4paper]{article}
\providecommand\ACNroot{../..}
\input{\ACNroot/templates/lecture-fa}

\ACNset{number}{14}
\ACNset{title}{مقدمه‌ای بر شبکه‌های نسل پنجم}
\ACNset{shorttitle}{مقدمهٔ نسل پنجم}
\ACNset{titleother}{5G Networks: Introduction}
\ACNset{part}{بخش چهارم -- سامانه و رویه‌های نسل پنجم (درس‌های ۱۴ و ۱۹ تا ۲۲)}
\ACNset{time}{۶ تا ۸ ساعت (۳ ساعت مطالعه، ۳ تا ۵ ساعت تمرین)}
\ACNset{level}{متوسط}
\ACNset{prereq}{درس‌های ۵ تا ۷ و ۹؛ مبانی \lr{OFDM}؛ ظرفیت شانون}
\ACNset{gen}{5}

\begin{document}
\ACNcover

% ---------------------------------------------------------------------
\section{اهداف یادگیری}
پس از مطالعهٔ این درس باید بتوانید:
\begin{objectives}
  \item سناریوهای استفاده و نیازمندی‌های کمینهٔ کلیدی \lr{IMT-2020} را \textbf{بیان کنید} و هر یک را به یک انتخاب طراحی
        نسل پنجم \textbf{مرتبط کنید}.
  \item عددشناسی مقیاس‌پذیر \acro{NR} را \textbf{توضیح دهید} و طول شکاف و نماد را برای هر فاصلهٔ زیرحامل \textbf{محاسبه کنید}.
  \item نرخ دادهٔ اوج تقریبی یک حامل \acro{NR} را با فرمول \acro{3GPP} \textbf{محاسبه کنید}.
  \item استقرارهای غیرمستقل (\lr{EN-DC}) و مستقل را \textbf{مقایسه کنید}.
  \item تأخیر، چگالی اتصال و ظرفیت ناحیه را برای استقرارهای داده‌شده \textbf{تخمین بزنید} و برآورده‌شدن یک نیازمندی را
        \textbf{داوری کنید}.
  \item هستهٔ نسل پنجم و موضوعات درس‌های بعدی را \textbf{طرح‌ریزی کنید}.
\end{objectives}

% ---------------------------------------------------------------------
\section{دانش پیش‌نیاز}
به معماری \acro{LTE} و پروتکل‌های رادیویی آن (درس‌های ۵ تا ۷)، \acro{EPC} (درس ۹)، ساختار شبکهٔ منابع \acro{OFDM}
(زیرحامل‌ها، نمادها، پیشوند چرخشی، بلوک‌های منبع) و فرمول ظرفیت شانون $C = B\log_2(1+\SNR)$
\cite{shannon1948mathematical} نیاز دارید.

\paragraph{واژه‌های کلیدی} واژگان این درس همراه با معادل انگلیسی آن‌ها:
\begin{keyterms}
\lr{IMT-2020} & \lr{IMT-2020} & نام \lr{ITU-R} برای خانوادهٔ واسط‌های رادیویی نسل پنجم و نیازمندی‌های آن. \\
سه سناریوی استفاده & \lr{eMBB / URLLC / mMTC} & پهن‌باند پیشرفته، فوق‌قابل‌اطمینان کم‌تأخیر، انبوه ماشینی. \\
رادیوی جدید & \lr{NR} & واسط رادیویی نسل پنجم \acro{3GPP} (از انتشار ۱۵ به بعد). \\
عددشناسی & \lr{Numerology} $\mu$ & فاصلهٔ زیرحامل $15\cdot2^\mu$ کیلوهرتز و طول شکاف متناظر. \\
بازه‌های بسامدی ۱ و ۲ & \lr{FR1 / FR2} & بازه‌های زیر ۷ گیگاهرتز و موج میلی‌متری. \\
بخش پهنای باند & \lr{Bandwidth part} & زیرمجموعه‌ای از یک حامل که برای یک \acro{UE} پیکربندی می‌شود. \\
اتصال دوگانهٔ \lr{E-UTRA--NR} & \lr{EN-DC} & گرهٔ اصلی \acro{LTE} و گرهٔ ثانویهٔ \acro{NR} با \acro{EPC} (\acro{NSA}). \\
مستقل / غیرمستقل & \lr{SA / NSA} & \acro{NR} با هستهٔ نسل پنجم / \acro{NR} لنگرشده روی \acro{LTE}. \\
هستهٔ نسل پنجم & \lr{5GC} & هستهٔ سرویس‌محور (\lr{AMF, SMF, UPF}، \ldots). \\
برش شبکه & \lr{Network slice} & شبکهٔ منطقی سرتاسری برای یک نوع خدمت. \\
\end{keyterms}

% ---------------------------------------------------------------------
\section{مقدمه و انگیزه}
\acro{LTE} برای یک کاربرد غالب طراحی شده بود: پهن‌باند موبایل برای گوشی‌های هوشمند. در میانهٔ دههٔ ۲۰۱۰ اپراتورها و
صنایع سه نوع تقاضای متفاوت را پیش‌بینی می‌کردند: ظرفیت بسیار بیشتر برای ویدئو و دسترسی بی‌سیم ثابت، کنترل بی‌سیم
ماشین‌ها با قابلیت اطمینان و تأخیری هم‌تراز شبکه‌های سیمی صنعتی، و اتصال برای تعداد بسیار زیادی حسگر ارزان و
باتری‌دار. \lr{ITU-R} این‌ها را به‌صورت سه سناریوی استفادهٔ \lr{IMT-2020} (\cref{fig:usage}) \cite{itu2015m2083} صورت‌بندی
کرد و نیازمندی‌های فنی کمینه‌ای برای هر واسط رادیویی نامزد تعیین کرد \cite{itu2017m2410}. پاسخ \acro{3GPP} سامانهٔ
نسل پنجم است: واسط رادیویی \acro{NR} \cite{ts38300,dahlman2020nr} و هستهٔ سرویس‌محور نسل پنجم \cite{ts23501}.

این درس بخش چهارم را می‌گشاید. نیازمندی‌ها، ایده‌های اصلی \acro{NR} که به آن‌ها پاسخ می‌دهند و مهاجرت از \acro{LTE} را مرور
می‌کند؛ شبکهٔ هسته، رویه‌ها، برش‌بندی و کیفیت خدمات در درس‌های ۱۸ تا ۲۲ می‌آیند، و شبکه‌سازی نرم‌افزارمحور و مجازی‌سازی،
فناوری‌هایی که هستهٔ نسل پنجم بر آن‌ها بنا شده، موضوع بخش سوم‌اند.

\begin{acnfigure}
  \acnfig{lecture-14/usage-scenarios}
  \caption{سه سناریوی استفادهٔ \lr{IMT-2020} با کاربردهای نمونه \cite{itu2015m2083}.}
  \label{fig:usage}
\end{acnfigure}

% ---------------------------------------------------------------------
\section{مفاهیم بنیادی}

\subsection{نیازمندی‌ها}
\begin{table}[htbp]
\centering\small
\caption{برخی نیازمندی‌های فنی کمینهٔ \lr{IMT-2020} (\lr{ITU-R M.2410}) \cite{itu2017m2410}.}\label{tab:req}
\rowcolors{2}{acnmist}{white}
\begin{tabularx}{\linewidth}{@{}l l Y@{}}
\toprule
\rowcolor{white}نیازمندی & مقدار & سناریو / توضیح \\
\midrule
نرخ دادهٔ اوج              & \lr{DL 20 Gbit/s, UL 10 Gbit/s}  & \acro{eMBB} \\
بازدهی طیفی اوج            & \lr{DL 30, UL 15 bit/s/Hz}       & \acro{eMBB} \\
نرخ دادهٔ تجربه‌شدهٔ کاربر  & \lr{DL 100 Mbit/s, UL 50 Mbit/s} & \acro{eMBB}، شهری متراکم (صدک پنجم) \\
ظرفیت ترافیک ناحیه         & \lr{10 Mbit/s/m$^2$}             & \acro{eMBB}، نقطهٔ داغ داخلی، فروسو \\
تأخیر صفحهٔ کاربر          & ۴ میلی‌ثانیه (\acro{eMBB})، ۱ میلی‌ثانیه (\acro{URLLC}) & یک‌طرفه، بدون بار \\
تأخیر صفحهٔ کنترل          & ۲۰ میلی‌ثانیه                   & از بیکار به فعال \\
چگالی اتصال                & $10^6$ دستگاه بر کیلومتر مربع    & \acro{mMTC} \\
قابلیت اطمینان             & $1-10^{-5}$ برای ۳۲ بایت در ۱ میلی‌ثانیه & \acro{URLLC} \\
تحرک                       & تا ۵۰۰ کیلومتر بر ساعت          & قطارهای سریع‌السیر \\
پهنای باند                 & دست‌کم ۱۰۰ مگاهرتز؛ تا ۱ گیگاهرتز بالای ۶ گیگاهرتز & -- \\
\bottomrule
\end{tabularx}
\end{table}

هیچ استقرار واحدی لازم نیست همهٔ مقادیر را هم‌زمان برآورده کند: هر نیازمندی در محیط آزمون خودش ارزیابی می‌شود.
\Cref{fig:capabilities} بهبودی را که \lr{ITU-R} برای هر قابلیت کلیدی نسبت به \lr{IMT-Advanced} (نسل چهارم) در نظر داشت
نشان می‌دهد.

\begin{acnfigure}
  \acnfig{lecture-14/capabilities}
  \caption{ضرایب بهبود قابلیت‌های کلیدی \lr{IMT-2020} نسبت به \lr{IMT-Advanced}، مطابق تصویر \lr{ITU-R M.2083}
           \cite{itu2015m2083}. تأخیر به‌صورت معکوس نشان داده شده است.}
  \label{fig:capabilities}
\end{acnfigure}

\subsection{اصول طراحی \lr{NR}}
\begin{itemize}
  \item \textbf{گسترهٔ وسیع طیف}: \lr{FR1} (۴۱۰ مگاهرتز تا ۷٫۱۲۵ گیگاهرتز) و \lr{FR2} (۲۴٫۲۵ تا ۵۲٫۶ گیگاهرتز، گسترش‌یافته در
        انتشارهای بعدی)، با حامل‌هایی تا ۱۰۰ مگاهرتز در \lr{FR1} و ۴۰۰ مگاهرتز در \lr{FR2} \cite{ts38300}.
  \item \textbf{عددشناسی مقیاس‌پذیر}: فاصلهٔ زیرحامل $15\cdot2^\mu$ کیلوهرتز (\cref{sec:num}).
  \item \textbf{طراحی فوق‌لاغر}: بدون سیگنال‌های مرجع همیشه‌روشن ویژهٔ سلول؛ سیگنال‌ها فقط در صورت نیاز ارسال می‌شوند که
        انرژی را صرفه‌جویی و تداخل را کم می‌کند.
  \item \textbf{عملکرد پرتومحور}: بلوک‌های سیگنال همگام‌سازی در پرتوهای مختلف جاروب می‌شوند؛ مدیریت پرتو جزء ذاتی است که
        در امواج میلی‌متری ضروری است.
  \item \textbf{زمان‌بندی انعطاف‌پذیر}: شکاف‌ها و شکاف‌های کوچک (۲، ۴ یا ۷ نماد)، \acro{TDD} پویا، بخش‌های پهنای باند.
  \item \textbf{کدگذاری کانال جدید}: \acro{LDPC} برای داده و کدهای قطبی برای کنترل \cite{ts38211}.
  \item \textbf{سازگاری رو به جلو}: منابع رزروشده تا قابلیت‌های آینده \acro{UE}های قدیمی را مختل نکنند.
\end{itemize}

\subsection{گزینه‌های استقرار}
\acro{3GPP} چند روش برای ترکیب \acro{LTE}، \acro{NR}، \acro{EPC} و \lr{5GC} تعریف کرد. دو روشی که در عمل مهم‌اند
(\cref{fig:options}):
\begin{itemize}
  \item \textbf{غیرمستقل (\acro{NSA})، خانوادهٔ گزینهٔ ۳ (\lr{EN-DC})}: \acro{UE} به یک \acro{eNB} \acro{LTE} (گرهٔ اصلی، صفحهٔ کنترل
        به \acro{MME}) و یک \lr{en-gNB} \acro{NR} (گرهٔ ثانویه که ظرفیت می‌افزاید) وصل است \cite{ts37340}. این روش \acro{EPC} را
        دوباره به کار می‌برد و نخستین راه‌اندازی‌های نسل پنجم در ۲۰۱۹ را بدون هستهٔ جدید ممکن کرد، اما هیچ‌یک از قابلیت‌های
        هستهٔ جدید (برش‌بندی، کیفیت خدمات جدید، معماری سرویس‌محور) را فراهم نمی‌کند.
  \item \textbf{مستقل (\acro{SA})، گزینهٔ ۲}: \acro{gNB} از طریق \lr{N2} (کنترل) و \lr{N3} (صفحهٔ کاربر) به \lr{5GC} وصل
        است. همهٔ قابلیت‌های نسل پنجم، از جمله تأخیر کم و برش‌بندی، به \acro{SA} نیاز دارند.
\end{itemize}

\begin{acnfigure}
  \resizebox{\linewidth}{!}{\acnfig{lecture-14/deployment-options}}
  \caption{\lr{EN-DC} غیرمستقل (گزینهٔ \lr{3x}) با \acro{EPC} و \acro{NR} مستقل (گزینهٔ ۲) با \lr{5GC}.}
  \label{fig:options}
\end{acnfigure}

% ---------------------------------------------------------------------
\section{توضیحات فنی تفصیلی}

\subsection{ساختار قاب}
\acro{NR} قاب ۱۰ میلی‌ثانیه‌ای و زیرقاب ۱ میلی‌ثانیه‌ای \acro{LTE} را حفظ می‌کند. هر شکاف ۱۴ نماد \acro{OFDM} دارد (پیشوند
چرخشی عادی)؛ با عددشناسی $\mu$ هر زیرقاب $2^\mu$ شکاف دارد (\cref{fig:numerology}). هر بلوک منبع ۱۲ زیرحامل در بسامد است؛
یک حامل حداکثر ۲۷۵ بلوک منبع دارد. چون ۱ میلی‌ثانیه شبکه‌ای ثابت برای همهٔ عددشناسی‌هاست، نمادهای عددشناسی‌های مختلف در
مرز زیرقاب هم‌تراز می‌شوند و این امکان آمیختن آن‌ها روی یک حامل را می‌دهد \cite{ts38211}.

\begin{acnfigure}
  \acnfig{lecture-14/numerology}
  \caption{تعداد شکاف‌ها در هر زیرقاب ۱ میلی‌ثانیه‌ای برای عددشناسی‌های $\mu = 0$ تا ۳.}
  \label{fig:numerology}
\end{acnfigure}

\subsection{دسترسی اولیه و پرتوها}
\acro{gNB} بلوک‌های سیگنال همگام‌سازی (\acro{SSB}: \acro{PSS}، \acro{SSS} و \acro{PBCH}) را در یک انفجار و هر یک در پرتوی
متفاوت ارسال می‌کند. \acro{UE} بهترین \acro{SSB} را تشخیص می‌دهد، اطلاعات سامانه را می‌خواند و دسترسی تصادفی را در فرصت
\acro{RACH} مرتبط با آن \acro{SSB} انجام می‌دهد؛ بدین ترتیب \acro{gNB} بهترین پرتو را می‌آموزد. در طول اتصال، پرتوها با
اندازه‌گیری‌ها و گزارش‌های \lr{CSI-RS} پالایش می‌شوند. در \lr{FR1} چند پرتوی پهن کافی است؛ در \lr{FR2} تا ۶۴ پرتوی باریک
تلفات مسیر بیشتر را با بهرهٔ آرایه جبران می‌کنند.

\subsection{از \lr{EPC} تا \lr{5GC}}
هستهٔ نسل پنجم مدیریت دسترسی و تحرک (\acro{AMF}) را از مدیریت نشست (\acro{SMF}) و صفحهٔ کاربر (\acro{UPF}) جدا می‌کند،
دادهٔ اشتراک را در \lr{UDM/UDR} ذخیره می‌کند، از طریق \acro{AUSF} احراز هویت می‌کند و سیاست را در \acro{PCF} کنترل می‌کند.
توابع شبکه خدمات خود را روی \lr{HTTP/2} عرضه می‌کنند و یکدیگر را از طریق \acro{NRF} کشف می‌کنند \cite{ts23501}. مدل کیفیت
خدمات به‌جای حامل‌های \acro{EPS} از جریان‌های کیفیت خدمات درون نشست‌های \acro{PDU} استفاده می‌کند، و برش‌های شبکه چند
شبکهٔ منطقی را روی یک زیرساخت ممکن می‌کنند. این موضوعات در درس‌های ۱۸ تا ۲۲ گسترش می‌یابند.

% ---------------------------------------------------------------------
\section{مدل‌های ریاضی}

\subsection{عددشناسی}\label{sec:num}
\paragraph{تعریف‌ها} برای $\mu = 0,1,2,3,4$:
\begin{equation}\label{eq:numerology}
  \Delta f = 15\cdot2^{\mu}\ \text{kHz}, \qquad
  T_{\mathrm{slot}} = \frac{1\ \text{ms}}{2^{\mu}}, \qquad
  \bar T_{\mathrm{sym}} = \frac{T_{\mathrm{slot}}}{14},
\end{equation}
که در آن $\bar T_{\mathrm{sym}}$ طول میانگین نماد همراه با پیشوند چرخشی است. پیشوند چرخشی عادی حدود $4.7/2^\mu$ میکروثانیه
است.
\paragraph{تفسیر} فاصلهٔ بیشتر نمادها و شکاف‌ها را کوتاه‌تر می‌کند (تأخیر کمتر) و نویز فاز و داپلر بیشتری را تحمل می‌کند (لازم در
\lr{FR2})، اما پیشوند چرخشی کوتاه‌ترش پخش تأخیر کمتری را تحمل می‌کند که با سلول‌های کوچک سازگار است. سلول‌های کلان
\lr{FR1} معمولاً از ۱۵ یا ۳۰ کیلوهرتز استفاده می‌کنند؛ \lr{FR2} از ۱۲۰ کیلوهرتز.

\subsection{نرخ دادهٔ اوج}\label{sec:peak}
\paragraph{فرمول} \acro{3GPP} بیشینهٔ نرخ دادهٔ یک حامل را چنین تقریب می‌زند \cite{ts38306}:
\begin{equation}\label{eq:peak}
  R \approx 10^{-6}\, v\, Q_m\, f\, R_{\max}\,
          \frac{N_{\mathrm{PRB}}\cdot 12}{T_s^{\mu}}\,\bigl(1-\mathrm{OH}\bigr)\ \text{Mbit/s},
  \qquad T_s^{\mu} = \frac{10^{-3}}{14\cdot2^{\mu}}\ \text{s},
\end{equation}
با $v$ لایهٔ \acro{MIMO}، $Q_m$ بیت در هر نماد مدولاسیون، ضریب مقیاس $f\le1$، بیشینهٔ نرخ کد $R_{\max} = 948/1024$،
$N_{\mathrm{PRB}}$ بلوک منبع و سربار \lr{OH} (۰٫۱۴ در فروسوی \lr{FR1}، ۰٫۱۸ در فروسوی \lr{FR2}، ۰٫۰۸ در فراسوی \lr{FR1} و
۰٫۱۰ در فراسوی \lr{FR2}). برای تجمیع حامل، نرخ‌های حامل‌ها جمع می‌شوند.
\paragraph{مثال} \lr{FR1}، ۱۰۰ مگاهرتز با ۳۰ کیلوهرتز ($N_{\mathrm{PRB}} = 273$)، ۴ لایه، \lr{256-QAM} ($Q_m = 8$)، $f = 1$:
$R \approx 2.34$ گیگابیت بر ثانیه (\cref{ex:peak}).

\subsection{تأخیر صفحهٔ کاربر}\label{sec:lat}
\paragraph{مدل} یک تخمین سادهٔ تأخیر رادیویی یک‌طرفه:
\begin{equation}\label{eq:latency}
  T_{\mathrm{UP}} \approx T_{\mathrm{align}} + T_{\mathrm{tx}} + T_{\mathrm{proc}}
  \approx \tfrac12 T_{\mathrm{slot}} + T_{\mathrm{slot}} + 2\,T_{\mathrm{slot}},
\end{equation}
که به‌طور میانگین نیم شکاف انتظار برای فرصت ارسال بعدی، یک شکاف ارسال و حدود دو شکاف پردازش در فرستنده و گیرنده را
می‌شمارد.
\paragraph{مثال} $\mu = 0$: ۳٫۵ میلی‌ثانیه؛ $\mu = 1$: ۱٫۷۵ میلی‌ثانیه؛ $\mu = 3$: ۰٫۴۴ میلی‌ثانیه. شکاف‌های کوچک و فراسوی
ازپیش‌پیکربندی‌شده (بدون اعطا) $T_{\mathrm{align}}$ و $T_{\mathrm{tx}}$ را بیشتر کوتاه می‌کنند.
\paragraph{محدودیت‌ها} مدل از درخواست‌های زمان‌بندی، ارسال‌های مجدد \acro{HARQ}، صف و شبکهٔ انتقال و هسته چشم می‌پوشد؛ نشان
می‌دهد چرا هدف ۱ میلی‌ثانیه‌ای \acro{URLLC} به عددشناسی‌های بزرگ‌تر یا شکاف‌های کوچک نیاز دارد.

\subsection{چگالی اتصال و ظرفیت ناحیه}\label{sec:density}
\paragraph{مدل‌ها} با چگالی دستگاه $\rho$ (دستگاه بر کیلومتر مربع) و شعاع سلول $r$، هر سلول به $N = \rho\pi r^2$ دستگاه خدمت
می‌دهد؛ اگر هر یک در هر دورهٔ $T$ یک بار به شبکه دسترسی یابد، نرخ دسترسی $N/T$ است. برای ظرفیت ناحیه، سلولی با پهنای باند
$B$ و بازدهی طیفی میانگین $\bar\eta$ که ناحیهٔ $A_{\mathrm{cell}}$ را پوشش می‌دهد این ظرفیت را عرضه می‌کند:
\begin{equation}\label{eq:area}
  C_{\mathrm{area}} = \frac{B\,\bar\eta}{A_{\mathrm{cell}}}.
\end{equation}
\paragraph{مثال‌ها} $\rho = 10^6$ بر کیلومتر مربع و $r = 500$ متر: $N = 7.9\times10^5$ دستگاه؛ یک گزارش در هر ۲ ساعت ۱۰۹
دسترسی در ثانیه در هر سلول می‌دهد (\cref{ex:density}). برای $C_{\mathrm{area}} = 10$ مگابیت بر ثانیه بر متر مربع با
$B = 400$ مگاهرتز و $\bar\eta = 5$ بیت بر ثانیه بر هرتز، هر سلول فقط می‌تواند $A_{\mathrm{cell}} = 200$ متر مربع را پوشش دهد:
ظرفیت ناحیه با متراکم‌سازی به دست می‌آید، نه فقط با بازدهی طیفی.

% ---------------------------------------------------------------------
\section{معماری شبکه و نمودارها}
\Cref{fig:options} دو معماری اصلی را خلاصه می‌کند. در \acro{NSA}، \acro{EPC} هسته باقی می‌ماند و گرهٔ \acro{NR} امتدادی از
شبکهٔ رادیویی \acro{LTE} است (اتصال صفحهٔ کنترل به هسته ندارد). در \acro{SA}، \acro{gNB} (که به واحدهای مرکزی و توزیع‌شده
تقسیم می‌شود، درس ۲۳) به \lr{5GC} وصل است؛ آن‌گاه \acro{LTE} و \acro{NR} برای تحرک و بازگشت صدا از طریق هسته (\lr{N26} میان
\acro{AMF} و \acro{MME}) همکاری می‌کنند. بیشتر اپراتورها با \acro{NSA} آغاز کردند و \acro{SA} را بعداً معرفی کردند، اغلب
نخست برای بخش‌های خاصی مانند دسترسی بی‌سیم ثابت یا برش‌های سازمانی.

% ---------------------------------------------------------------------
\section{رویه‌های گام‌به‌گام}

\begin{procedurebox}{افزودن گرهٔ ثانویهٔ \lr{NR} در \lr{EN-DC}}\label{proc:endc}
\begin{steps}
  \item \acro{UE} به \acro{eNB} \acro{LTE} (گرهٔ اصلی) وصل است و اندازه‌گیری‌های \acro{NR} پیکربندی‌شده توسط آن را گزارش می‌دهد.
  \item \acro{eNB} تصمیم می‌گیرد یک شاخهٔ \acro{NR} بیفزاید و \lr{SgNB Addition Request} را از طریق \lr{X2} به \lr{en-gNB}
        می‌فرستد.
  \item \lr{en-gNB} منابع را تخصیص می‌دهد و با پیکربندی رادیویی \acro{NR} پاسخ می‌دهد؛ \acro{eNB} آن را در
        \lr{RRC Connection Reconfiguration} هدایت می‌کند.
  \item \acro{UE} دسترسی تصادفی را روی سلول \acro{NR} انجام می‌دهد و تأیید می‌کند؛ \acro{eNB} به \acro{MME} اطلاع می‌دهد
        (اصلاح \lr{E-RAB}) تا \lr{S-GW} بتواند صفحهٔ کاربر را مستقیم به \lr{en-gNB} بفرستد (گزینهٔ \lr{3x}).
  \item اکنون داده روی \acro{NR} (و اگر حامل تقسیم‌شده باشد، روی \acro{LTE}) جریان دارد؛ تحرک و کنترل نزد \acro{LTE}
        می‌مانند.
\end{steps}
\end{procedurebox}

\begin{procedurebox}{دسترسی اولیه در \lr{NR} مستقل}\label{proc:access}
\begin{steps}
  \item شبکهٔ \acro{SSB} را جست‌وجو کن، \lr{PSS/SSS} را تشخیص بده و شناسهٔ سلول و بهترین پرتوی \acro{SSB} را تعیین کن.
  \item \acro{MIB} را روی \acro{PBCH} و \lr{SIB1} (که با آن زمان‌بندی می‌شود) را بخوان تا پیکربندی \acro{RACH} و پارامترهای
        سلول به دست آید.
  \item یک پیش‌آهنگ را در فرصت \acro{RACH} بهترین \acro{SSB} ارسال کن؛ پاسخ دسترسی تصادفی را دریافت کن (درس ۸).
  \item برقراری \acro{RRC} را کامل کن و \lr{Registration Request} را به \acro{AMF} بفرست (درس ۲۰).
\end{steps}
\end{procedurebox}

% ---------------------------------------------------------------------
\section{مثال‌های حل‌شده}

\begin{examplebox}{نرخ اوج یک حامل \lr{FR1}}\label{ex:peak}
\textbf{مسئله.} \cref{eq:peak} را برای ۱۰۰ مگاهرتز با ۳۰ کیلوهرتز، ۴ لایه، \lr{256-QAM}، $f = 1$ در فروسو ارزیابی کنید.

\textbf{حل.} $T_s^{1} = 10^{-3}/28 = 35.7$ میکروثانیه. $N_{\mathrm{PRB}}\cdot12/T_s = 273\times12/35.7\times10^{-6} = 9.17\times10^{7}$
عنصر منبع در ثانیه. $R = 4\times8\times0.926\times9.17\times10^{7}\times0.86 = 2.34\times10^{9}$ بیت بر ثانیه. برای رسیدن به ۲۰ گیگابیت
بر ثانیه، \acro{NR} به پهنای باندهای \lr{FR2}، لایه‌های بیشتر و تجمیع حامل نیاز دارد.
\end{examplebox}

\begin{examplebox}{تأخیر و عددشناسی}\label{ex:lat}
\textbf{مسئله.} با مدل \cref{eq:latency} کدام عددشناسی‌ها هدف ۱ میلی‌ثانیه‌ای \acro{URLLC} را برآورده می‌کنند؟

\textbf{حل.} $T_{\mathrm{UP}} = 3.5\,T_{\mathrm{slot}} \le 1$ میلی‌ثانیه مستلزم $T_{\mathrm{slot}} \le 0.286$ میلی‌ثانیه است، یعنی
$\mu \ge 2$ (۰٫۸۷۵ میلی‌ثانیه با ۶۰ کیلوهرتز). با ۳۰ کیلوهرتز، شکاف‌های کوچک ۲ نمادی جملات ارسال و هم‌ترازی را کاهش می‌دهند
و آن‌ها نیز می‌توانند هدف را برآورده کنند.
\end{examplebox}

\begin{examplebox}{بار دسترسی \lr{mMTC}}\label{ex:density}
\textbf{مسئله.} سلولی به شعاع ۵۰۰ متر در ناحیه‌ای با $10^6$ دستگاه بر کیلومتر مربع؛ هر دستگاه هر دو ساعت یک گزارش می‌فرستد.
نرخ دسترسی را محاسبه و با ظرفیت معمول \acro{RACH} در \acro{LTE}، در حد چند صد تلاش در ثانیه برای هر سلول، مقایسه کنید.

\textbf{حل.} $N = 10^6\times\pi\times0.25 = 7.85\times10^5$؛ نرخ $= 7.85\times10^5/7200 = 109$\,s$^{-1}$ به‌طور میانگین -- قابل
مدیریت، اما گزارش‌دهی همگام (همهٔ دستگاه‌ها سر ساعت) فاجعه‌بار خواهد بود (درس ۸). بنابراین \acro{mMTC} به زمان‌بندی‌های
تصادفی‌شده، ارسال زودهنگام داده و قابلیت‌های \lr{NB-IoT/LTE-M} تکیه دارد، نه فقط به ظرفیت خام.
\end{examplebox}

% ---------------------------------------------------------------------
\section{مقایسهٔ فناوری‌ها}
\begin{table}[htbp]
\centering\small
\caption{مقایسهٔ \acro{LTE} و \acro{NR}.}\label{tab:cmp}
\rowcolors{2}{acnmist}{white}
\begin{tabularx}{\linewidth}{@{}l Y Y@{}}
\toprule
\rowcolor{white} & \acro{LTE} (نسل ۴) & \acro{NR} (نسل ۵) \\
\midrule
طیف                 & زیر ۶ گیگاهرتز                  & \lr{FR1} و \lr{FR2} (موج میلی‌متری) \\
بیشینهٔ پهنای باند حامل & ۲۰ مگاهرتز (تجمیع تا ۳۲ حامل) & ۱۰۰ مگاهرتز (\lr{FR1})، ۴۰۰ مگاهرتز (\lr{FR2}) \\
فاصلهٔ زیرحامل       & ۱۵ کیلوهرتز (ثابت)               & $15\cdot2^\mu$ کیلوهرتز \\
واحد زمان‌بندی       & زیرقاب ۱ میلی‌ثانیه‌ای           & شکاف ($1/2^\mu$ میلی‌ثانیه)، شکاف کوچک \\
سیگنال‌های مرجع      & \acro{CRS} همیشه‌روشن           & در صورت نیاز (فوق‌لاغر) \\
کدگذاری کانال        & توربو (داده)، کانولوشنی (کنترل)  & \acro{LDPC} (داده)، قطبی (کنترل) \\
شکل‌دهی پرتو         & در انتشارهای بعدی افزوده شد      & ذاتی، از جمله در دسترسی اولیه \\
هسته                 & \acro{EPC}                       & \acro{EPC} (\acro{NSA}) یا \lr{5GC} (\acro{SA}) \\
واحد کیفیت خدمات     & حامل \acro{EPS}                   & جریان کیفیت خدمات \\
\bottomrule
\end{tabularx}
\end{table}

\paragraph{تکامل، نه جایگزینی} \acro{LTE} به تکامل ادامه می‌دهد (\lr{LTE-M}، \lr{NB-IoT} و اشتراک پویای طیف با \acro{NR} روی
یک حامل)، و \acro{NR} بسیاری از مفاهیم \acro{LTE} را دوباره به کار می‌برد: \acro{OFDM}، \acro{HARQ}، پشتهٔ لایهٔ ۲ و منطق
اتصال/ثبت‌نام و حامل/جریان در هسته.

% ---------------------------------------------------------------------
\section{نکته‌های کلیدی}
\begin{takeawaybox}
\begin{itemize}
  \item \lr{IMT-2020} سه سناریو (\lr{eMBB, URLLC, mMTC}) و نیازمندی‌های کمینه‌ای مانند اوج ۲۰ گیگابیت بر ثانیه، تأخیر ۱
        میلی‌ثانیه‌ای \acro{URLLC} و $10^6$ دستگاه بر کیلومتر مربع تعریف می‌کند.
  \item \acro{NR}: طیف گسترده (\lr{FR1}، \lr{FR2})، عددشناسی $15\cdot2^\mu$ کیلوهرتز، شکاف‌های $1/2^\mu$ میلی‌ثانیه‌ای، لاغر و
        پرتومحور.
  \item نرخ اوج $\propto v\,Q_m\,R_{\max}\,N_{\mathrm{PRB}}\cdot12/T_s\,(1-\mathrm{OH})$؛ ۱۰۰ مگاهرتز در \lr{FR1} با ۴ لایه حدود ۲٫۳
        گیگابیت بر ثانیه می‌دهد.
  \item \acro{NSA} (\lr{EN-DC}) از \acro{EPC} استفاده می‌کند؛ \acro{SA} با \lr{5GC} برش‌بندی، کیفیت خدمات جدید و تأخیر کم را
        ممکن می‌کند.
  \item ظرفیت ناحیه از متراکم‌سازی، تأخیر از عددشناسی و شکاف‌های کوچک، و چگالی از طراحی دسترسی به دست می‌آید.
\end{itemize}
\end{takeawaybox}

% ---------------------------------------------------------------------
\section{تمرین‌ها}

\begin{exercisebox}{1}{کدام سناریو؟}\label{exr:scenario}
هر خدمت را به \acro{eMBB}، \acro{URLLC} یا \acro{mMTC} نسبت دهید و سخت‌گیرانه‌ترین نیازمندی آن را نام ببرید: (الف) پخش
ویدئوی \lr{8K}؛ (ب) کنترل حرکت در کارخانه؛ (ج) کنتورهای آب؛ (د) جراحی از راه دور؛ (هـ) دسترسی بی‌سیم ثابت به خانه‌ها.
\end{exercisebox}

\begin{exercisebox}{1}{چرا نخست \lr{NSA}؟}\label{exr:nsa}
دو دلیل بیاورید که چرا بیشتر اپراتورها نسل پنجم را به‌صورت \acro{NSA} راه‌اندازی کردند، و یک قابلیت که تا \acro{SA}
نمی‌توانستند ارائه دهند.
\end{exercisebox}

\begin{exercisebox}{2}{نرخ اوج \lr{FR2}}\label{exr:fr2}
نرخ اوج فروسوی یک حامل ۴۰۰ مگاهرتزی \lr{FR2} با ۱۲۰ کیلوهرتز ($N_{\mathrm{PRB}} = 264$) با ۲ لایه و \lr{64-QAM} را محاسبه
کنید. برای ۲۰ گیگابیت بر ثانیه چند حامل از این نوع لازم است؟
\end{exercisebox}

\begin{exercisebox}{2}{پیشوند چرخشی و اندازهٔ سلول}\label{exr:cp}
پیشوند چرخشی عادی حدود $4.7/2^\mu$ میکروثانیه است. آن را برای $\mu = 0$ تا ۳ و اختلاف طول مسیری که می‌تواند جذب کند محاسبه
کنید. کدام عددشناسی‌ها برای یک سلول کلان روستایی با پخش تأخیر ۲ میکروثانیه مناسب‌اند؟
\end{exercisebox}

\begin{exercisebox}{3}{بازدهی طیفی اوج}\label{exr:se}
با \cref{eq:peak} بازدهی طیفی اوج یک حامل ۱۰۰ مگاهرتزی \lr{FR1} با ۳۰ کیلوهرتز، ۸ لایهٔ فروسو و \lr{256-QAM} را محاسبه کنید.
آیا نیازمندی ۳۰ بیت بر ثانیه بر هرتز \lr{IMT-2020} را برآورده می‌کند؟ کدام فرض‌ها این را مقداری نظری می‌کنند؟
\end{exercisebox}

\begin{exercisebox}{3}{متراکم‌سازی}\label{exr:dense}
یک نقطهٔ داغ داخلی باید ۱۰ مگابیت بر ثانیه بر متر مربع فراهم کند. ناحیهٔ سلول لازم را برای (الف) ۴۰۰ مگاهرتز با ۵ بیت بر ثانیه
بر هرتز و (ب) ۱۰۰ مگاهرتز با ۷ بیت بر ثانیه بر هرتز مقایسه کنید. پیامدها برای بک‌هال و شبکهٔ هسته را بررسی کنید.
\end{exercisebox}

\subsection*{پاسخ‌ها}
\addcontentsline{toc}{subsection}{پاسخ‌ها}

\begin{solutionbox}{\ref{exr:scenario}}
(الف) \acro{eMBB} -- نرخ دادهٔ کاربر؛ (ب) \acro{URLLC} -- تأخیر و قابلیت اطمینان؛ (ج) \acro{mMTC} -- چگالی اتصال و عمر باتری؛
(د) \acro{URLLC} (همراه با ویدئوی \acro{eMBB}) -- قابلیت اطمینان و تأخیر؛ (هـ) \acro{eMBB} -- ظرفیت در هر ناحیه و هر کاربر.
\end{solutionbox}

\begin{solutionbox}{\ref{exr:nsa}}
\acro{NSA} از \acro{EPC} موجود (بدون هستهٔ جدید، ورود سریع‌تر به بازار) و پوشش \acro{LTE} برای کنترل و تحرک استفاده می‌کرد،
در حالی که \acro{NR} ظرفیت می‌افزود؛ دستگاه‌ها و طیف اولیه نیز برای آن تنظیم شده بودند. قابلیت‌هایی که به \lr{5GC} نیاز دارند --
برش‌بندی شبکه، جریان‌های کیفیت خدمات، عرضهٔ سرویس‌محور، صفحه‌های کاربر لبه با تأخیر کم و پوشش فقط \acro{NR} -- به \acro{SA}
نیاز داشتند.
\end{solutionbox}

\begin{solutionbox}{\ref{exr:fr2}}
$T_s^{3} = 10^{-3}/112 = 8.93$ میکروثانیه؛ $R = 2\times6\times0.926\times264\times12/8.93\times10^{-6}\times0.82 = 3.23$ گیگابیت بر
ثانیه. برای ۲۰ گیگابیت بر ثانیه حدود $20/3.23 = 6.2$، یعنی ۷ حامل -- یا لایه‌های بیشتر و مدولاسیون مرتبهٔ بالاتر.
\end{solutionbox}

\begin{solutionbox}{\ref{exr:cp}}
$\mu = 0$: ۴٫۶۹ میکروثانیه (۱۴۰۶ متر)؛ $\mu = 1$: ۲٫۳۴ میکروثانیه (۷۰۳ متر)؛ $\mu = 2$: ۱٫۱۷ میکروثانیه (۳۵۲ متر)؛ $\mu = 3$:
۰٫۵۹ میکروثانیه (۱۷۶ متر)، با $d = c\,T_{\mathrm{CP}}$. با پخش تأخیر ۲ میکروثانیه، $\mu = 0$ (و به‌طور مرزی $\mu = 1$) مناسب‌اند؛
۶۰ کیلوهرتز به پیشوند چرخشی گسترش‌یافته نیاز دارد و ۱۲۰ کیلوهرتز برای سلول‌های کوچک \lr{FR2} است.
\end{solutionbox}

\begin{solutionbox}{\ref{exr:se}}
با ۸ لایه نرخ دو برابر می‌شود و به $4.67$ گیگابیت بر ثانیه می‌رسد، یعنی $46.7$ بیت بر ثانیه بر هرتز روی ۱۰۰ مگاهرتز -- بالاتر از
۳۰. این مقدار ۸ لایهٔ جداپذیر فضایی، \lr{256-QAM} با بالاترین نرخ کد (\acro{SINR} بسیار زیاد)، ضریب سربار ثابت و تخصیص کامل به
یک \acro{UE} را فرض می‌کند؛ سامانه‌های واقعی به کسری از آن می‌رسند.
\end{solutionbox}

\begin{solutionbox}{\ref{exr:dense}}
(الف) $400\times10^6\times5 = 2$ گیگابیت بر ثانیه برای هر سلول، پس $A = 2\times10^9/10^7 = 200$ متر مربع؛ (ب) $0.7$ گیگابیت بر
ثانیه، پس ۷۰ متر مربع. در هر دو حالت سلول‌ها به اندازهٔ یک اتاق‌اند، نه یک خیابان. هر سلول به بک‌هال چندگیگابیتی (فیبر)
نیاز دارد و هسته باید ترافیک $10$ ترابیت بر ثانیه در هر کیلومتر مربع از چنین نقاط داغی را به‌صورت محلی مدیریت کند -- استدلالی
به سود صفحه‌های کاربر لبه (درس ۲۳).
\end{solutionbox}

% ---------------------------------------------------------------------
\section{مطالعهٔ بیشتر}
\lr{ITU-R M.2083} \cite{itu2015m2083} چشم‌انداز و \lr{M.2410} \cite{itu2017m2410} نیازمندی‌ها را بیان می‌کنند؛ سناریوهای خود
\acro{3GPP} در \lr{TR 38.913} \cite{tr38913} و نیازمندی‌های خدمت در \lr{TS 22.261} \cite{ts22261} آمده‌اند. \acro{NR} در
\lr{TS 38.300} \cite{ts38300} و \lr{TS 38.211} \cite{ts38211}، \lr{EN-DC} در \lr{TS 37.340} \cite{ts37340} و سامانهٔ نسل
پنجم در \lr{TS 23.501} \cite{ts23501} توصیف شده‌اند. دالمن و همکاران \cite{dahlman2020nr} کتاب درسی استاندارد \acro{NR} است.

\begin{latin}
\printbibliography[heading=subbibliography,title={References}]
\end{latin}
\end{document}
